% Scheletro di lavoro, non testo definitivo.
% Le coordinate indicate sono piste di lettura da verificare sul testo greco prima della stesura.
% Metodo: testo originale, traduzione, ricostruzione argomentativa, analisi critica.
\chapter{Essere, non essere e differenza nel Sofista}
\label{cap:platone:sofista}

% Problema: In che modo la definizione del sofista rivela problemi logici e ontologici?
% Fonti da verificare: Sofista 216a–236d.
\section{La ricerca del sofista}

\subsection{La divisione come metodo}

\subsection{Immagine, imitazione e apparente sapere}

% Problema: Come è possibile dire ciò che non è senza contraddizione?
% Fonti da verificare: Sofista 237a–242b.
\section{Il problema del non essere}

\subsection{L'eredità eleatica}

\subsection{L'aporia del discorso falso}

% Problema: Quali difficoltà emergono dalle dottrine sull'essere?
% Fonti da verificare: Sofista 242b–251a.
\section{Essere, movimento e quiete}

\subsection{La battaglia fra giganti e dèi}

\subsection{Essere, movimento e quiete}

% Problema: In quale senso il non essere può essere inteso come differenza?
% Fonti da verificare: Sofista 251a–264b.
\section{La comunanza dei generi}

\subsection{I generi sommi e la loro partecipazione}

\subsection{Il diverso, il discorso falso e la definizione del sofista}
